\subsection{Further consequences of quotient-cardinality shifts}

The same argument applies to a ring with finitely many ideals of each
finite index. Write $a_n$ for their number, and assume that
$L(s)=\sum_{n\ge1}a_n n^{-s}$ has finite abscissa of absolute convergence
$\sigma_a$. Set $L_t(s)=\sum_{n\ge1}a_n(n+t)^{-s}$. The following statements
also hold for complex coefficients $a_n$ whenever the stated absolute
convergence hypothesis holds.

\begin{proposition}\label{prop:quotient-shift-remainder}
For each real $t\ge0$, $L_t$ and $L$ have the same abscissa of absolute
convergence. If $L$ is meromorphic on $\mathbb C$, then for $|t|<1$ and
$N\ge1$,
\[
 L_t(s)=\sum_{j=0}^{N-1}\frac{(-1)^j(s)_j}{j!}t^j L(s+j)+H_{N,t}(s),
 \qquad \Re s>\sigma_a-N,
\]
where $H_{N,t}$ is holomorphic on the displayed half-plane. If the only
pole of $L$ is a simple pole at $1$, then $L_t-L$ is entire, so the pole
and residue are unchanged.
\end{proposition}
\begin{proof}
For each fixed real $\sigma$, the ratios
$(n+t)^{-\sigma}/n^{-\sigma}$ are bounded above and below by positive
constants, since $1\le(n+t)/n\le1+t$. Comparison of absolute values proves
the abscissa assertion.

Taylor's formula for $(1+z)^{-s}$ gives
\[
 (n+t)^{-s}=\sum_{j<N}\frac{(-1)^j(s)_j}{j!}t^j n^{-s-j}
       +t^N n^{-s-N}R_N(s,t/n).
\]
For $s$ in a compact set and $|t|\le r<1$, the integral remainder is
uniformly bounded: its integrand is a polynomial factor in $s$ times
$(1+u t/n)^{-s-N}$, $0\le u\le1$, whose base stays in the right half-plane
away from zero. Thus the remainder series converges absolutely and
locally uniformly when $\Re s>\sigma_a-N$. It defines $H_{N,t}$ and the
continuation. The definitions for different $N$ agree on their initial
half-plane and hence agree throughout their common domain.

For $j\ge1$, the only possible pole of $L(s+j)$ is at $s=1-j$. The factor
$(s)_j=s(s+1)\cdots(s+j-1)$ vanishes there and cancels that simple pole.
Every term with $j\ge1$ is therefore entire. Taking $N$ arbitrarily large
proves that $L_t-L$ is entire for $|t|<1$. For any fixed real $t\ge0$,
choose an integer $A>t$ and remove the finitely many terms with $n<A$.
The same Taylor argument applies to the remaining series, uniformly
because $t/n\le t/A<1$. The finite removed terms are entire functions
of $s$, and the truncated unshifted series has the same pole and
residue as $L$. Thus the same pole cancellations prove the claim for
every nonnegative shift. This cancellation at the displaced poles
is needed in addition to checking the residue at $1$.
\end{proof}

If $N(I)$ is a positive multiplicative ordinary ideal norm, the shifted
weight $w(I)=N(I)+1$ obeys
\[
 w(I)w(J)-w(IJ)=N(I)+N(J)>0.
\]
Its quotient-counting series cannot have a normalized absolutely
convergent ordinary ideal Euler product: for real $s\to+\infty$ the
series tends to zero, since every shifted base is at least $2$, whereas
such an Euler product tends to its constant coefficient $1$. For
$\mathbb Z$ it is $\zeta(s)-1$. The ordinary sheet for a number ring is
its Dedekind ideal zeta function, directly from the definition by ideal
norms. Finiteness of quotients alone does not imply multiplicativity of
those norms, as the order $\mathbb Z[2i]$ and its ideal $(2,2i)$ show:
the indices of that ideal and its square are $2$ and $8$.

\begin{proposition}\label{prop:quotient-log-jets}
For the quotient-counting path $L_t(s)=\zeta(s,1+t)$, and $\Re s>1$, put
\[
 \alpha_r(n)=(-1)^{\omega(n)}
       \frac{\prod_{p\mid n}(p^r-1)}{n^r}\quad(r\ge1),
 \qquad A_r(s)=\sum_{n\ge1}\alpha_r(n)n^{-s}.
\]
The empty product gives $\alpha_r(1)=1$. Then
\[
 \left.\partial_t\log L_t(s)\right|_{t=0}=-sA_1(s),
 \qquad
 \left.\partial_t^2\log L_t(s)\right|_{t=0}
       =s(s+1)A_2(s)-s^2A_1(s)^2.
\]
These are meromorphic identities after continuation wherever the
logarithmic derivatives are defined.
\end{proposition}
\begin{proof}
In the region of absolute convergence,
\[
 \frac{\zeta(s+r)}{\zeta(s)}
 =\prod_p\frac{1-p^{-s}}{1-p^{-s-r}}.
\]
Writing $u=p^{-s}$, each factor expands as
$1-\sum_{k\ge1}(p^r-1)p^{-rk}u^k$. Multiplying the absolutely convergent
factors gives $A_r(s)=\zeta(s+r)/\zeta(s)$. The first two ordinary
logarithmic-derivative formulas, together with
$\partial_t^jL_t(s)=(-1)^j(s)_jL_t(s+j)$, give the displayed identities.
The identity theorem gives their meromorphic continuations; no global
choice of a logarithm across zeros is asserted.
\end{proof}

\begin{proposition}[The Bernoulli endpoint correction]
\label{prop:bernoulli-endpoint}
For real $t\ge0$, the quotient-sheet defect
$D_t(s)=\zeta(s)-\zeta(s,1+t)$ is entire and satisfies
\[
 D_t(-m)=\frac{B_{m+1}(1+t)-B_{m+1}(1)}{m+1}\quad(m\ge0),
 \qquad D_t'(0)=-\log\Gamma(1+t).
\]
In particular $D_t(0)=t$. Replacing $B_1(1)=1/2$ by the conventional
Bernoulli number $B_1(0)=-1/2$ gives an erroneous extra $1$ at $m=0$.
\end{proposition}
\begin{proof}
Both Hurwitz functions have their sole simple pole at $1$ with residue
$1$, so their difference is entire. Apply the classical identities
$\zeta(-m,a)=-B_{m+1}(a)/(m+1)$ and
$\zeta'(0,a)=\log\Gamma(a)-\tfrac12\log(2\pi)$, at $a=1+t$ and $a=1$;
see \citet[\S25.11(i), equations 25.11.14 and 25.11.18]{DLMFShift}.
Subtracting gives both formulas, and $B_1(x)=x-1/2$ gives the endpoint.
\end{proof}
